% things to ask serra:
% why does div thm hold over piecewise manifolds, what is even the normal defined as 
% why do the coefficients of the differential forms have to be smooth (infty)
% i don't think we ever used it, just needed it to ensure dalpha is also a smooth form
% what is the deal with "positively oriented" parametrisation, and how exactly is the normal vector to the manifold boundary defined

\chapter{Ordinary Differential Equations}

This part is a bit self-contained and independent of the rest of the course. The goal is to finish off ordinary differential equations with this for good.

\section{Constant coefficient linear equations}

We start with the simplest type of ODE that actually has a complete answer to existence and uniqueness of solutions.

\begin{definition}[Constant coefficient linear ODE]\label{def:const-lin-ode}
   A differential equation for $y: I \to \mathbb{R}$ of class $C^{m}$ (sometimes the codomain is just $\mathbb{C}$) where $I \subset \mathbb{R}$ is an interval, is constant coefficient and linear if it is of the form
    \begin{align*}
        \sum_{i=0}^{m} a_{i}y^{(i)}(t) = f(t) \quad \forall t \in I,
    \end{align*}
    where $f: I \to \mathbb{R}$ is class $C^{0}$, the coefficients $a_{i}$ are in $\mathbb{R}$ and $a_{m} = 1$. We call $m$ the order of the equation.
\end{definition}

\begin{definition}[Characteristic polynomial of constant linear ODE]\label{def:char-poly}
    The characteristic polynomial of constant linear ODE $\sum_{i=0}^{m} a_{i}y^{(i)}(t) = f(t)$ is defined as $p(X) = \sum_{i=0}^{m} a_{i}X^{i}$. Sometimes we summarise the differential equation as the linear operator $L = p\left( \frac{d}{dt} \right) = \sum_{i=0}^{m} a_{i}\left( \frac{d}{dt} \right)^{i}$, i.e. a linear homomorphism $C^{m}(I) \to C^{m}(I)$. Because then the ODE can be written compactly as $Ly = f$.
\end{definition}

\begin{definition}[Homogeneous ODE]\label{def:homogen-ode}
    The differential equation $L z = 0$ for $z: I \to \mathbb{C}$ is called the homogeneous ODE. The space of solutions to this equation $V = \text{Ker}(L)$ forms a vector space over $\mathbb{C}$.
\end{definition}

The goal for today is to completely characterise this space $V$. Being good at linear algebra, we know that all there really is to know is the basis for this space. So that's what we will be looking for. 

We start with the Ansatz $z(t) = \text{exp}(\alpha t)$ for $\alpha \in \mathbb{C}$ and notice this satisfies $\left( \frac{d}{dt} \right)^{k} z(t) = \alpha^{k} z(t)$. Hence
\begin{align*}
    Lz = p \left( \frac{d}{dt} \right) z = \left( \sum_{i=0}^{m} a_{i} \alpha^{i} \right)  z(t) = p(\alpha ) z(t) = 0
\end{align*}
if $\alpha$ is root of $p$. 

So for a typical polynomial $p$ with $m$ \textit{distinct} roots $\left\{ \alpha_{1}, \dots, \alpha_{m} \right\}$, we discover that $\dim V \geq m$ (show as an exercise that the solutions for different $\alpha_{i}$ are linearly independent).

However, what happens when there are roots of $p$ that have multiplicity? Notice that for $k \geq 2$,
\begin{align*}
    \left( \frac{d}{dt} - \alpha  \right)^{k} (t^{k-1} \text{exp}(\alpha t)) &= \left( \frac{d}{dt} - \alpha  \right)^{k-1}\left( \frac{d}{dt} - \alpha \right)(t^{k-1} \text{exp}(\alpha t)) \\ 
    &= \left( \frac{d}{dt} - \alpha  \right)^{k-1} ((k-1) t^{k-2} \text{exp}(\alpha t) - \alpha t^{k-1} \text{exp}(\alpha t) + t^{k-1}\alpha \cdot \text{exp}(\alpha t)),
\end{align*}
the terms cancel and we can apply induction. Hence $\left( \frac{d}{dt} - \alpha  \right)^{k}(t^{k-1} \text{exp}(\alpha t)) = 0$ for $k \geq 2$ (if $k=0,1$ anyway). 

\begin{lemma}[Poly-exp form]\label{lem:polyexp}
    Consider a differential equation with characteristic polynomial $p(X) = \prod_{j=1}^{k}(X-\alpha_{j})^{m_{j}}$, which has degree $m = \sum_{j=1}^{k} m_{j}$. Then for a fixed $1 \leq j \leq k$,
    \begin{align*}
        p \left( \frac{d}{dt} \right) (t^{k} \text{exp}(\alpha_{j} t)) = 0 \quad \forall 0 \leq k \leq m_{j}-1.
    \end{align*}
    If as an exercise, one proves these functions are all linearly independent, we can conclude that $\dim V \geq \text{deg}(p)$ in \textit{any} case.
\end{lemma}
\begin{proof}
    Fix a $j$ and a fitting $k$. Factorise the characteristic polynomial as
    \begin{align*}
        p \left( \frac{d}{dt} \right) = \left( \frac{d}{dt} - \alpha_{j} \right)^{m_{j}} \prod_{\ell \neq j} \left( \frac{d}{dt} - \alpha_{\ell} \right)^{m_{\ell}}.
    \end{align*}
    Hence
    \begin{align*}
        p \left( \frac{d}{dt} \right) \left( t^{k} \text{exp}(\alpha_{j} t) \right) = \left( \frac{d}{dt} - \alpha_{j} \right)^{m_{j}} \left( t^{k} \text{exp}(\alpha_{j} t) \right) \cdot (\cdots) = 0,
    \end{align*}
    as we found above. 
\end{proof}

The rest of the work today will go into showing that $\dim V \leq \text{deg}(p)$ to complete the theory of linear constant coefficient ODEs.

\begin{lemma}[Grönwell inequality]\label{lem:gronwall}
    Let $u \in C^{1}([a,b])$ such that $u'(t) \leq \beta(t) u(t)\; \forall t \in [a,b]$, where $\beta \in C^{0}([a,b])$. Then 
    \begin{align*}
        u(t) \leq u(a) \cdot \text{exp}\left( \int_{a}^{t} \beta (s) \, ds \right) \quad \forall t \in [a,b].
    \end{align*}
\end{lemma}
\begin{proof}
    Define $v(t) = \text{exp}\left( \int_{a}^{t} \beta (s) \, ds \right) \in C^{1}([a,b])$. By FTC, this satisfies $v'(t) = \beta (t)v(t)$. Take $\phi(t) := \frac{u(t)}{v(t)}$. Using the assumption,
    \begin{align*}
        \phi'(t) = \frac{u'(t)v(t) - u(t)v'(t)}{v(t)^{2}} \leq \frac{\beta (t) u(t)v(t) - u(t) \beta (t) v(t)}{v(t)^{2}} = 0.
    \end{align*}
    Hence $\phi(t) \leq \phi(a) \; \forall t \in [a,b]$. But when you plug in the definition for $\phi$, you get $\frac{u(t)}{v(t)} \leq \frac{u(a)}{1}$, which is just the statement.
\end{proof}

\begin{lemma}[Zero in, zero out]\label{lem:zero-in-out}
    Consider a differential equation for $w: I \to \mathbb{C}$ with characteristic polynomial $p(X) = \prod_{j=1}^{k}(X-\alpha_{j})^{m_{j}}$ and let $t_{0} \in I$ and the \textit{initial value problem} (IVP)
    \begin{align*}
        \begin{cases}
            Lw(t) = 0 \quad \forall t \in I \\ 
            w^{(j)}(t_{0}) = 0 \quad \forall 0 \leq j \leq m-1
        \end{cases}
    \end{align*}
    Then $w \equiv  0$ is the unique solution.
\end{lemma}
\begin{proof}
    Consider $u(t) = \sum_{i=0}^{m-1} (w^{(i)}(t))^{2}$. Then
\begin{align*}
    u'(t) = 2w^{(m-1)}(t) \cdot w^{(m)}(t) + \sum_{i=0}^{m-2} 2 w^{(i)}(t) \cdot w^{(i+1)}(t).
\end{align*}
Because $Lw = 0$, $ w^{(m)}(t) = - \sum_{i=0}^{m-1} a_{i}w^{(i)}(t)$ and therefore
\begin{align*}
    \left| w^{(m)}(t) \right|
    \leq \sum_{i=0}^{m-1} |a_i|\, |w^{(i)}(t)|
    \leq C \sum_{i=0}^{m-1} \left| w^{(i)}(t) \right|,
\end{align*}
with $C = \max_{i} |a_i|$. Hence, by $2ab \leq a^{2}+b^{2}$,
\begin{align*}
    u'(t)
    &\leq \sum_{i=0}^{m-2} \left( (w^{(i)}(t))^{2} + (w^{(i+1)}(t))^{2} \right)
        + 2 \left| w^{(m-1)}(t) \right| \left| w^{(m)}(t) \right| \\
    &\leq 2u(t)
        + 2C \left| w^{(m-1)}(t) \right|
        \sum_{i=0}^{m-1} \left| w^{(i)}(t) \right| \\
    &\leq 2u(t)
        + C \left( (w^{(m-1)}(t))^{2}
        + \left( \sum_{i=0}^{m-1} |w^{(i)}(t)| \right)^{2} \right).
\end{align*}
Using $\left( \sum_{i=0}^{m-1} |w^{(i)}(t)| \right)^{2} \leq m \sum_{i=0}^{m-1} (w^{(i)}(t))^{2} = mu(t)$, we obtain
\begin{align*}
    u'(t)
    \leq 2u(t) + C(u(t) + mu(t))
    = (2 + C(m+1))u(t).
\end{align*}
Applying Grönwall \ref{lem:gronwall}, we get
\begin{align*}
    0 \leq u(t)
    \leq u(t_{0}) \exp(\overline{C}(t-t_{0}))
\end{align*}
for $\overline{C} = 2 + C(m+1)$. Since $u(t_{0}) = 0$, we conclude that $u(t)=0$, and hence $w(t)=0$.
\end{proof}

\begin{prop}[Initial value problem]\label{prop:init-val-prob}
    Consider a differential equation for $z: I \to \mathbb{C}$ with characteristic polynomial $p(X) = \prod_{j=1}^{k}(X-\alpha_{j})^{m_{j}}$ and let $t_{0} \in I$. Consider the initial value problem 
    \begin{align*}
        \begin{cases}
            Lz(t) = 0 \quad \forall t \in I \\ 
            z^{(j)}(t_{0}) = w_{j} \quad \forall 0 \leq j \leq m-1
        \end{cases}
    \end{align*}
    for $w_{j} \in \mathbb{C}$. It has a \textit{unique} solution.
\end{prop}
\begin{proof}
    Consider the homomorphism $T: V \to \mathbb{C}^{m}: z \mapsto (z(t_{0}), z^{(1)}(t_{0}), \dots, z^{(m-1)}(t_{0}))$. By Lemma \ref{lem:zero-in-out}, $\text{Ker}(T) = 0$. Thus $T$ is an injective homomorphism from $V$ to $\mathbb{C}^{m}$. But that means $\dim V \leq \dim \mathbb{C}^{m} = m$.
    
    By Lemma \ref{lem:polyexp}, $\dim V = m$ and $T$ is an isomorphism assigning each IVP a unique solution.
\end{proof}

